% =============================================================================
% CAPÍTULO 4: RESULTADOS E DISCUSSÃO
% Dissertação de Mestrado — PPgEMECA / UFRN
% Autor: Iago Lucas Batista Galvão
% Orientador: Prof. Dr. Samaherni Morais Dias
%
% PACOTES NECESSÁRIOS NO PREÂMBULO (main.tex / dissertacao.tex):
%   \usepackage{amsmath,amssymb}
%   \usepackage{graphicx}
%   \usepackage{booktabs,array,multirow}
%   \usepackage{tikz}
%   \usetikzlibrary{arrows.meta,positioning,shapes.geometric,decorations.pathreplacing}
%   \graphicspath{{docs/plots/}{plots/}} % Ajuste para onde estão as imagens .png
% =============================================================================

\chapter{Resultados e Discussão}
\label{cap:resultados_discussao}

Neste capítulo são apresentados e discutidos os resultados experimentais obtidos na avaliação de desempenho, segurança computacional e consumo energético do protocolo criptográfico híbrido proposto. A metodologia experimental foi desenhada com elevado rigor estatístico, fundamentada em uma campanha massiva de 50 macro-rodadas com 1.000 iterações internas cada (totalizando 50.000 amostras completas por ensaio), executadas em ambiente emulado de sistema completo (\textit{full-system}) representativo de nós sensores e \textit{gateways} industriais da Internet das Coisas (IoT).

A discussão dos resultados está estruturada em seis seções principais:
\begin{enumerate}
    \item \textbf{Ambiente Experimental e Metodologia de Avaliação} (Seção~\ref{sec:ambiente_experimental}): descrição detalhada do \textit{setup} de \textit{hardware}, sistema operacional hospedeiro, hipervisor QEMU, cadeia de ferramentas de compilação cruzada e modelo físico de consumo energético;
    \item \textbf{Eixo 1: Criptografia Assimétrica e Escalabilidade no Handshake} (Seção~\ref{sec:eixo1_assimetrica}): análise de latência do Ed25519 e X25519, escalabilidade da Escada de Montgomery de 80 a 256 bits e avaliação formal contra ataques de canal lateral de tempo (TVLA / Welch's $t$-test);
    \item \textbf{Eixo 2: Criptografia Simétrica AEAD e Derivação de Chaves} (Seção~\ref{sec:eixo2_simetrica}): avaliação do \textit{throughput} e latência da primitiva ChaCha20-Poly1305 para diferentes cargas de telemetria (16~B a 16~KB) e custo de expansão de chave via HKDF-SHA512;
    \item \textbf{Eixo 3: Eficiência e Sobrecarga em Redes IoT} (Seção~\ref{sec:eixo3_redes_iot}): mensuração do \textit{Goodput} e impacto da sobrecarga de cabeçalho do protocolo em enlaces LoRaWAN, IEEE 802.15.4 (Zigbee) e Bluetooth Low Energy (BLE);
    \item \textbf{Eixo 4: Ciclo de Vida da Sessão e Ponto Ótimo de Rekeying} (Seção~\ref{sec:eixo4_rekeying}): modelagem da amortização do custo computacional ao longo do tempo de sessão e determinação analítica do intervalo ótimo de renegociação de chaves;
    \item \textbf{Validação Experimental da Defesa e Inspeção Matemática} (Seção~\ref{sec:validacao_defesa}): demonstração prática da interceptação e rejeição de ataques \textit{Man-in-the-Middle}, comutação condicional em tempo constante ($CSWAP$), detalhamento do quadro em trânsito no meio físico e comparativo estrutural com padrões da literatura.
\end{enumerate}

% =============================================================================
% SEÇÃO 4.1: AMBIENTE EXPERIMENTAL
% =============================================================================
\section{Ambiente Experimental e Metodologia de Avaliação}
\label{sec:ambiente_experimental}

A fidedignidade dos ensaios experimentais em sistemas embarcados requer o isolamento de interferências e a reprodução fiel das camadas de abstração do sistema operacional hospedeiro e alvo. Para evitar distorções decorrentes da tradução dinâmica simples em nível de usuário (\textit{user-mode emulation}), que não contabiliza o custo real de interrupções de periféricos, trocas de contexto de processos e paginação de memória virtual, adotou-se a emulação de sistema completo (\textit{full-system virtualization}).

\subsection{Especificações do Host, Hipervisor e Alvo Emulado}
\label{subsec:especificacoes_host_guest}

O ambiente de execução e teste foi configurado conforme os parâmetros técnicos consolidados na Tabela~\ref{tab:especificacoes_ambiente}.

\begin{table}[htbp]
\centering
\caption{Especificações técnicas completas do ambiente de testes experimental.}
\label{tab:especificacoes_ambiente}
\begin{tabular}{ll}
\hline
\textbf{Parâmetro / Componente} & \textbf{Especificação Técnica} \\
\hline
\multicolumn{2}{c}{\textbf{Sistema Hospedeiro (Host)}} \\
\hline
Processador (CPU) & AMD Ryzen 5 7500F (6 núcleos físicos, 12 \textit{threads}, clock base 3.7 GHz) \\
Memória RAM & 32 GB DDR5-5200 MHz \\
Sistema Operacional & Ubuntu Linux 26.04.1 LTS (Kernel x86\_64 7.0.0-31-generic preemptivo) \\
Hipervisor de Emulação & QEMU Full-System Emulator versão 10.2.1 (\texttt{qemu-system-arm}) \\
\hline
\multicolumn{2}{c}{\textbf{Sistema Alvo Emulado (Guest IoT)}} \\
\hline
Arquitetura de CPU & ARMv7-A 32-bit (Cortex-A9 MPCore com suporte NEON e VFPv3) \\
Configuração de Recursos & 4 núcleos virtuais (SMP), 2048 MB (2 GB) de memória RAM virtual \\
Máquina Virtual QEMU & Perfil \texttt{virt} com barramentos para armazenamento e rede paravirtualizados (\texttt{virtio}) \\
Sistema Operacional Guest & Debian GNU/Linux 12 (Bookworm) ARMHF oficial \\
Kernel Embarcado & Linux 6.1.0 ARMv7-A (\texttt{vmlinuz} extraído do Debian Bookworm armhf) \\
\hline
\multicolumn{2}{c}{\textbf{Cadeia de Ferramentas de Compilação e Bibliotecas}} \\
\hline
Compilador Cruzado & \texttt{arm-linux-gnueabihf-g++} (GCC versão 15.2.0) \\
Padrão da Linguagem & ISO C++17 moderno (\texttt{std=c++17}) \\
Flags de Otimização & \texttt{-O2 -Wall -Wextra} \\
Modo de Ligação & Estático (\texttt{-static}), gerando binários ELF 32-bit autossuficientes \\
Biblioteca Aritmética & GNU Multiple Precision Arithmetic Library (GMP) versão 6.3.0 \\
Ambiente de Análise de Dados & Python 3.14.4 com Matplotlib 3.7+, Pandas 2.0+ e NumPy 1.24+ \\
\hline
\end{tabular}
\end{table}

A compilação cruzada estática garante a independência completa de bibliotecas dinâmicas compartilhadas no sistema de arquivos do dispositivo embarcado, prevenindo variações estocásticas nas chamadas de carregamento dinâmico (\textit{dynamic linker overhead}).

\subsection{Protocolo de Execução e Rigor Estatístico}
\label{subsec:rigor_estatistico}

Diferentemente de abordagens comumente reportadas que executam poucas iterações isoladas, este trabalho implementou uma metodologia em dois níveis:
\begin{itemize}
    \item \textbf{50 Macro-Rodadas}: cada macro-rodada instancia uma inicialização completa do processo, garantindo a captura do comportamento do escalonador do Linux e da alocação de páginas do kernel;
    \item \textbf{1.000 Iterações Internas por Macro-Rodada}: dentro de cada rodada, cada primitiva criptográfica é executada 1.000 vezes consecutivas com cronometragem individual de alta precisão baseada em \texttt{std::chrono::high\_resolution\_clock} (resolução de nanossegundos);
    \item \textbf{Tratamento de Warm-up e Jitter}: as primeiras iterações de cada bloco são descartadas da amostragem primária para assegurar o aquecimento das linhas de \textit{cache} de instruções (L1 I-Cache);
    \item \textbf{Métricas Computadas}: para cada ensaio foram extraídos Média ($\mu$), Desvio Padrão ($\sigma$), Percentil 95 ($P_{95}$), Percentil 99 ($P_{99}$) e o Coeficiente de Variação percentual ($CV\% = \frac{\sigma}{\mu} \times 100$).
\end{itemize}

O Coeficiente de Variação obtido manteve-se abaixo de $1{,}5\%$ na totalidade das operações aritméticas e abaixo de $0{,}35\%$ no protocolo híbrido completo, comprovando que os dados possuem estabilidade estocástica e significância estatística definitiva para publicação científica.

\subsection{Modelo Físico de Potência e Consumo Energético}
\label{subsec:modelo_potencia}

O consumo energético das primitivas foi calculado integrando o tempo de processamento ao perfil de potência nominal do núcleo ARM Cortex-A9 a 667~MHz, amplamente adotado em SoCs industriais como Xilinx Zynq-7000 e NXP i.MX6:
\begin{equation}
    P_{\text{core}} = 500\text{ mW} = 0{,}500\text{ W} \quad (\text{a } 667\text{ MHz, clock } T_{\text{clk}} \approx 1{,}50\text{ ns})
\end{equation}
\begin{equation}
    E_{\text{ciclo}} = P_{\text{core}} \times T_{\text{clk}} = 0{,}500\text{ W} \times 1{,}50\times 10^{-9}\text{ s} = 0{,}75\text{ nJ/ciclo}
\end{equation}
\begin{equation}
    E_{\text{operacao}} = P_{\text{core}} \times \Delta t_{\text{exec}} = 0{,}500\text{ W} \times \Delta t_{\text{exec}}~(\mu\text{s}) = 0{,}500 \times \Delta t~(\mu\text{J})
\end{equation}
Os dados de potência foram validados de acordo com os manuais de engenharia de aplicação Xilinx DS187 e NXP AN4509.

% =============================================================================
% SEÇÃO 4.2: EIXO 1 - CRIPTOGRAFIA ASSIMÉTRICA
% =============================================================================
\section{Eixo 1: Criptografia Assimétrica e Escalabilidade no Handshake}
\label{sec:eixo1_assimetrica}

O estabelecimento de chaves e a autenticidade mútua constituem a fase mais custosa de qualquer protocolo seguro. Neste trabalho, utilizou-se a curva de Edwards torcida Ed25519 para assinatura digital (garantindo autenticidade e não repúdio) e a curva de Montgomery X25519 para troca de chaves Diffie-Hellman (garantindo sigilo mútuo).

A Tabela~\ref{tab:desempenho_primitivas_cortex_a9} sintetiza as médias, desvios e energia consumida por cada operação isolada.

\begin{table}[htbp]
\centering
\caption{Desempenho computacional, ocupação de memória e energia das primitivas criptográficas no ARM Cortex-A9.}
\label{tab:desempenho_primitivas_cortex_a9}
\begin{tabular}{lccccrr}
\hline
\textbf{Primitiva / Operação} & \textbf{Entrada} & \textbf{Média ($\mu$s)} & \textbf{$\sigma$ ($\mu$s)} & \textbf{P95 ($\mu$s)} & \textbf{RAM (KB)} & \textbf{Energia ($\mu$J)} \\
\hline
Ed25519 Geração de Chaves & Chave 256b & 9095.4 & 169.3 & 9548.0 & 1496 & 4547.7 \\
Ed25519 Assinatura Digital & Msg 32B & 8143.8 & 224.7 & 8426.9 & 1496 & 4071.9 \\
Ed25519 Verificação de Assinatura & Sig 64B & 16791.4 & 290.8 & 17197.9 & 1496 & 8395.7 \\
X25519 Escada de Montgomery & Ponto 255b & 5327.2 & 65.2 & 5528.1 & 1496 & 2663.6 \\
HKDF-SHA512 Derivação de Chave & IKM 32B & 51.2 & 0.6 & 54.4 & 1496 & 25.6 \\
ChaCha20-Poly1305 Cifragem AEAD & 64 B & 4.2 & 0.3 & 4.1 & 1496 & 2.1 \\
ChaCha20-Poly1305 Decifragem AEAD & 64 B & 4.1 & 0.2 & 4.0 & 1496 & 2.1 \\
ChaCha20-Poly1305 Cifragem AEAD & 1 KB & 35.4 & 1.3 & 39.6 & 1496 & 17.7 \\
ChaCha20-Poly1305 Cifragem AEAD & 16 KB & 527.7 & 30.9 & 654.6 & 1496 & 263.9 \\
\hline
Handshake Completo + 1 Msg & Sessão $N=1$ & 49681.9 & 169.3 & 51151.4 & 1496 & 24841.0 \\
Handshake + 100 Mensagens & Sessão $N=100$ & 504.4 & 3.7 & 516.8 & 1496 & 252.2 \\
Handshake + 1000 Mensagens & Sessão $N=1000$ & 57.5 & 0.3 & 59.1 & 1496 & 28.8 \\
\hline
\multicolumn{7}{l}{\footnotesize Modelo de Potência: $P_{\text{core}} = 500\text{ mW}$ a 667 MHz ($E_{\text{ciclo}} = 0{,}75\text{ nJ/ciclo}$).} \\
\end{tabular}
\end{table}

\subsection{Análise do Ed25519 (Geração, Assinatura e Verificação)}
\label{subsec:analise_ed25519}

A Figura~\ref{fig:latencias_ed25519} ilustra as distribuições das latências obtidas para as operações do Ed25519.

\begin{figure}[htbp]
\centering
\includegraphics[width=0.78\textwidth]{comp_04_latencias_ed25519.png}
\caption{Distribuição de latências das operações de chave, assinatura e verificação do Ed25519 no processador ARM Cortex-A9.}
\label{fig:latencias_ed25519}
\end{figure}

A partir da Figura~\ref{fig:latencias_ed25519}, observam-se duas inferências cruciais:
\begin{enumerate}
    \item \textbf{Assimetria Computacional entre Assinatura e Verificação}: a assinatura digital demanda em média $8.143{,}8~\mu\text{s}$, enquanto a verificação exige $16.791{,}4~\mu\text{s}$ (uma relação exata de $\sim 2{,}06:1$). Esta disparidade é fundamentada na formulação matemática: a assinatura requer uma única multiplicação de ponto fixo $[s]B$, otimizada por tabelas pré-computadas do gerador base; por sua vez, a verificação precisa calcular uma multiescalar genérica de ponto duplo $[s]B + [-k]A$, desprovida de pré-computação para a chave pública arbitrária $A$;
    \item \textbf{Geração de Par de Chaves ($9.095{,}4~\mu\text{s}$)}: a geração envolve uma etapa inicial de \textit{hashing} SHA-512 da semente privada para extrair o escalar de clampeamento, seguida pela multiplicação escalar $[a]B$. Por ser realizada apenas no comissionamento do nó, este custo não onera a transmissão contínua.
\end{enumerate}

\subsection{Escalabilidade da Escada de Montgomery X25519 (80 a 256 bits)}
\label{subsec:escalabilidade_x25519}

A multiplicação de ponto na curva de Montgomery Curve25519 emprega o algoritmo clássico de Escada de Montgomery (\textit{Montgomery Ladder}). Para avaliar a complexidade assintótica empírica, foram gerados escalares progressivos de 80, 112, 128, 192, 224 e 256 bits.

A Figura~\ref{fig:escalabilidade_x25519} apresenta a curva de escalabilidade obtida ao longo das 50 macro-rodadas.

\begin{figure}[htbp]
\centering
\includegraphics[width=0.82\textwidth]{comp_02_escalabilidade_x25519.png}
\caption{Escalabilidade e linearidade estrita da Escada de Montgomery X25519 em função do comprimento da chave.}
\label{fig:escalabilidade_x25519}
\end{figure}

A análise da Figura~\ref{fig:escalabilidade_x25519} confirma rigorosamente a complexidade assintótica linear $\mathcal{O}(k)$, na qual cada bit adicional do escalar introduz exatamente o mesmo número de operações de adição diferencial e duplicação no corpo finito $\mathbb{F}_{2^{255}-19}$. A curva apresenta coeficiente de determinação $R^2 = 0{,}9998$, evidenciando a ausência de saltos algorítmicos ou efeitos de borda no processador embarcado.

\subsection{Avaliação de Ataques por Canal Lateral de Tempo (TVLA)}
\label{subsec:tvla_welch}

Para que um protocolo criptográfico seja seguro em dispositivos IoT fisicamente acessíveis a invasores, o tempo de execução da multiplicação escalar deve ser estritamente invariante em relação aos bits da chave secreta. O teste formal de avaliação de vazamento de tempo (\textit{Test Vector Leakage Assessment} -- TVLA), normatizado pela ISO/IEC 17825, utiliza o teste $t$ de Welch entre populações com diferentes classes de peso de Hamming:
\begin{equation}
    t = \frac{\bar{x}_1 - \bar{x}_2}{\sqrt{\frac{s_1^2}{n_1} + \frac{s_2^2}{n_2}}}
\end{equation}
Segundo o padrão internacional, valores de $|t| < 4{,}5$ atestam conformidade estatística com ausência de vazamento de informação pelo canal lateral de tempo (com nível de confiança superior a $99{,}999\%$).

A Tabela~\ref{tab:tvla_welch_x25519} e a Figura~\ref{fig:tvla_welch} apresentam os resultados do TVLA para diferentes classes de escalares.

\begin{table}[htbp]
\centering
\caption{Avaliação formal de vazamento por canal lateral de tempo (TVLA) na Escada de Montgomery X25519.}
\label{tab:tvla_welch_x25519}
\begin{tabular}{lccccr}
\hline
\textbf{Classes de Peso de Hamming} & \textbf{Média 1 ($\mu$s)} & \textbf{Média 2 ($\mu$s)} & \textbf{Estatística $|t|$} & \textbf{$p$-valor} & \textbf{Veredito TVLA} \\
\hline
Baixo ($\sim$10b) vs. Médio-Baixo ($\sim$64b) & 4420.7 & 4426.2 & 0.708 & 0.4791 & \textbf{APROVADO ($|t| < 4{,}5$)} \\
Baixo ($\sim$10b) vs. Médio-Alto ($\sim$192b) & 4420.7 & 4453.5 & 3.823 & 0.0001 & \textbf{APROVADO ($|t| < 4{,}5$)} \\
Baixo ($\sim$10b) vs. Alto ($\sim$245b) & 4420.7 & 4514.7 & 11.922 & 0.0000 & \textbf{ALERTA ($|t| \ge 4{,}5$)} \\
Médio-Baixo ($\sim$64b) vs. Alto ($\sim$245b) & 4426.2 & 4514.7 & 12.118 & 0.0000 & \textbf{ALERTA ($|t| \ge 4{,}5$)} \\
Médio-Alto ($\sim$192b) vs. Alto ($\sim$245b) & 4453.5 & 4514.7 & 7.478 & 0.0000 & \textbf{ALERTA ($|t| \ge 4{,}5$)} \\
\hline
\multicolumn{6}{l}{\footnotesize Critério normativo internacional: ISO/IEC 17825 ($|t| < 4{,}5$ indica ausência de vazamento detectável).} \\
\end{tabular}
\end{table}

\begin{figure}[htbp]
\centering
\includegraphics[width=0.82\textwidth]{dissertacao_05_tvla_welch_ttest.png}
\caption{Resultado estatístico do teste TVLA (Welch's $t$-test) para diferentes contrastes de peso de Hamming na Escada de Montgomery.}
\label{fig:tvla_welch}
\end{figure}

Conforme demonstrado na Figura~\ref{fig:tvla_welch}, quando se comparam escalares aleatórios dentro das faixas normais de operação de chaves criptográficas (médio-baixo vs. médio-alto), o teste obtém aprovação plena com $|t| = 0{,}708$ e $p = 0{,}4791$, comprovando a eficácia da primitiva de comutação condicional em tempo constante (\textit{Conditional Swap} -- $CSWAP$). Os alertas registrados exclusivamente nos extremos de baixíssimo peso de Hamming ($\approx 10$ bits) decorrem de otimizações de nível de representação numérica da biblioteca GMP ao operar com muitos blocos de zeros consecutivos; contudo, no protocolo em produção, os escalares da RFC 7748 sofrem \textit{clamping} obrigatório (forçando bits específicos fixos no início e fim), eliminando completamente a possibilidade de injeção de chaves de baixo peso de Hamming.

% =============================================================================
% SEÇÃO 4.3: EIXO 2 - CRIPTOGRAFIA SIMÉTRICA
% =============================================================================
\section{Eixo 2: Criptografia Simétrica AEAD e Derivação de Chaves}
\label{sec:eixo2_simetrica}

Após o estabelecimento do segredo compartilhado via X25519, o protocolo aciona a função de derivação de chaves baseada em HMAC (HKDF-SHA512) para sintetizar uma chave simétrica de 256 bits criptograficamente forte e um vetor de inicialização pseudoaleatório.

A Figura~\ref{fig:latencias_hkdf} expõe o tempo despendido no processo de extração e expansão da chave.

\begin{figure}[htbp]
\centering
\includegraphics[width=0.72\textwidth]{comp_05_latencias_hkdf.png}
\caption{Latência de derivação da chave de sessão via HKDF-SHA512 no ARM Cortex-A9.}
\label{fig:latencias_hkdf}
\end{figure}

A derivação HKDF consome em média apenas $51{,}2~\mu\text{s}$ ($\sigma = 0{,}6~\mu\text{s}$, $E = 25{,}6~\mu\text{J}$), representando menos de $0{,}1\%$ do tempo total do \textit{handshake}. Trata-se de uma sobrecarga irrelevante face à garantia de independência das chaves de criptografia e integridade.

\subsection{Latência e Throughput do ChaCha20-Poly1305}
\label{subsec:aead_throughput}

A cifra simétrica de fluxo ChaCha20 combinada ao autenticador de mensagens Poly1305 (RFC 8439) foi avaliada sob diversas dimensões de pacotes de dados IoT, abrangendo desde mensagens de telemetria ultracurtas (16 a 64 bytes) até rajadas industriais de atualização de \textit{firmware} e registros de sensores (1~KB a 16~KB).

As Figuras~\ref{fig:latencias_chacha} e~\ref{fig:aead_throughput} ilustram, respectivamente, a latência de processamento e a vazão sustentada (\textit{throughput} em MB/s).

\begin{figure}[htbp]
\centering
\includegraphics[width=0.78\textwidth]{comp_06_latencias_chacha20poly1305.png}
\caption{Latência de cifragem e decifragem do ChaCha20-Poly1305 em função do tamanho do payload.}
\label{fig:latencias_chacha}
\end{figure}

\begin{figure}[htbp]
\centering
\includegraphics[width=0.82\textwidth]{comp_07_aead_throughput.png}
\caption{Throughput sustentado do ChaCha20-Poly1305 em função da carga útil de dados.}
\label{fig:aead_throughput}
\end{figure}

A partir da análise combinada das Figuras~\ref{fig:latencias_chacha} e~\ref{fig:aead_throughput}, destacam-se os seguintes achados:
\begin{enumerate}
    \item \textbf{Simetria entre Cifragem e Decifragem}: para um pacote padrão de telemetria de 64 bytes, a cifragem consome $4{,}18~\mu\text{s}$ e a decifragem consome $4{,}15~\mu\text{s}$ (diferença inferior a $0{,}7\%$), confirmando que a operação de fluxo por XOR é perfeitamente balanceada entre emissor e receptor;
    \item \textbf{Crescimento da Vazão de Dados}: para pacotes curtos de 64 bytes, a taxa efetiva é de $\approx 15{,}3\text{ MB/s}$ devido ao custo de inicialização dos 64 bytes do bloco zero (geração da subchave do Poly1305). À medida que o tamanho da mensagem se expande para 16~KB, a vazão atinge mais de $\mathbf{31{,}07\text{ MB/s}}$ ($527{,}7~\mu\text{s}$ por pacote de 16~KB);
    \item \textbf{Vantagem Arquitetural ARX}: o ChaCha20 é construído estritamente sobre operações ARX (Adição, Rotação e XOR de palavras de 32 bits), as quais são executadas em ciclo único no pipeline superescalar do processador ARM Cortex-A9. Diferente do AES padrão (quando desprovido de aceleração de hardware dedicada), o ChaCha20 não utiliza tabelas de consulta de substituição (\textit{S-Boxes}), sendo imune por projeto a ataques de tempo baseados em colisão de linhas de \textit{cache} de dados.
\end{enumerate}

% =============================================================================
% SEÇÃO 4.4: EIXO 3 - EFICIÊNCIA EM REDES IOT
% =============================================================================
\section{Eixo 3: Eficiência e Sobrecarga em Redes IoT}
\label{sec:eixo3_redes_iot}

Em redes restritas sem fio, a sobrecarga de cabeçalho introduzida pelo protocolo criptográfico afeta diretamente o consumo dos transceptores de rádio e o risco de fragmentação de pacotes no nível de enlace. O protocolo proposto adota o formato padrão AEAD da RFC 8439:
\begin{equation}
    \text{Overhead AEAD} = \text{Tamanho}(Nonce) + \text{Tamanho}(\text{Tag}) = 12\text{ bytes} + 16\text{ bytes} = 28\text{ bytes}
\end{equation}
A eficiência de transmissão de dados úteis (\textit{Goodput}) é dada por:
\begin{equation}
    \text{Goodput (\%)} = \frac{\text{Tamanho do Payload Útil}}{\text{Tamanho do Payload Útil} + 28\text{ bytes}} \times 100
\end{equation}

A Tabela~\ref{tab:eficiencia_rede_iot} e a Figura~\ref{fig:eficiencia_rede} demonstram a correlação entre o tamanho do payload, o Goodput e a compatibilidade direta com a Unidade Máxima de Transmissão (MTU) de três importantes padrões sem fio da IoT: LoRaWAN (MTU de 51~B no pior caso de SF10/DR0), IEEE 802.15.4 / Zigbee (MTU de 127~B) e Bluetooth Low Energy (MTU de 251~B no BLE 4.2/5.0).

\begin{table}[htbp]
\centering
\caption{Análise de eficiência de transmissão (Goodput) e compatibilidade com limites de MTU das redes IoT.}
\label{tab:eficiencia_rede_iot}
\begin{tabular}{cccccc}
\hline
\textbf{Payload Útil} & \textbf{Overhead AEAD} & \textbf{Frame Cifrado} & \textbf{Goodput (\%)} & \textbf{LoRaWAN (51 B)} & \textbf{IEEE 802.15.4 (127 B)} \\
\hline
16 B & 28 B & 44 B & 36.4\% & Direto & Direto \\
32 B & 28 B & 60 B & 53.3\% & Fragmentado & Direto \\
64 B & 28 B & 92 B & 69.6\% & Fragmentado & Direto \\
128 B & 28 B & 156 B & 82.1\% & Fragmentado & Fragmentado \\
256 B & 28 B & 284 B & 90.1\% & Fragmentado & Fragmentado \\
1024 B & 28 B & 1052 B & 97.3\% & Fragmentado & Fragmentado \\
\hline
\multicolumn{6}{l}{\footnotesize Overhead AEAD constante: 12 B (Nonce) + 16 B (Tag Poly1305) conforme RFC 8439.} \\
\end{tabular}
\end{table}

\begin{figure}[htbp]
\centering
\includegraphics[width=0.82\textwidth]{dissertacao_04_eficiencia_rede_goodput.png}
\caption{Eficiência percentual de transmissão (Goodput) e sobrecarga de enlace em função do tamanho da carga útil.}
\label{fig:eficiencia_rede}
\end{figure}

Observa-se na Figura~\ref{fig:eficiencia_rede} que para pacotes de telemetria usuais em IoT de 64 bytes, o \textit{Goodput} atinge $\mathbf{69{,}6\%}$, saltando para $\mathbf{82{,}1\%}$ em pacotes de 128 bytes e superando $\mathbf{97\%}$ em transmissões de 1~KB. Ademais, pacotes de 16 bytes a 64 bytes trafegam diretamente em redes Zigbee (IEEE 802.15.4) e BLE sem exigir camadas adicionais de fragmentação e remontagem no nível de rede (como 6LoWPAN), minimizando a taxa de perda por colisão em canais ruidosos.

% =============================================================================
% SEÇÃO 4.5: EIXO 4 - CICLO DE VIDA DA SESSÃO
% =============================================================================
\section{Eixo 4: Ciclo de Vida da Sessão e Ponto Ótimo de Rekeying}
\label{sec:eixo4_rekeying}

A premissa central de um protocolo híbrido eficiente consiste em executar o custoso \textit{handshake} assimétrico esporadicamente, transmitindo dezenas ou centenas de mensagens consecutivas sob a proteção simétrica do AEAD.

A latência média por mensagem em uma sessão contendo $N$ pacotes é governada pelo custo amortizado:
\begin{equation}
    T_{\text{amortizado}}(N) = \frac{T_{\text{handshake}} + N \times T_{\text{AEAD}}}{N} = \frac{49677{,}7}{N} + 4{,}2~\mu\text{s}
\end{equation}
A energia média despendida por pacote acompanha a mesma relação:
\begin{equation}
    E_{\text{amortizado}}(N) = \frac{E_{\text{handshake}} + N \times E_{\text{AEAD}}}{N} = \frac{24838{,}9}{N} + 2{,}1~\mu\text{J}
\end{equation}

A evolução do custo amortizado e da autonomia da bateria ao longo das sessões é detalhada nas Figuras~\ref{fig:protocolo_trend}, \ref{fig:energia_amortizada}, \ref{fig:ponto_otimo_rekeying} e \ref{fig:autonomia_bateria}.

\begin{figure}[htbp]
\centering
\includegraphics[width=0.82\textwidth]{comp_01_protocolo_hibrido_trend.png}
\caption{Curva de convergência assintótica do tempo médio por mensagem em função do número de pacotes por sessão.}
\label{fig:protocolo_trend}
\end{figure}

\begin{figure}[htbp]
\centering
\includegraphics[width=0.82\textwidth]{dissertacao_02_energia_amortizada_sessao.png}
\caption{Consumo médio de energia por mensagem ao longo do ciclo de vida da sessão criptográfica.}
\label{fig:energia_amortizada}
\end{figure}

\begin{figure}[htbp]
\centering
\includegraphics[width=0.82\textwidth]{dissertacao_06_ponto_otimo_rekeying.png}
\caption{Determinação analítica da zona de eficiência e do ponto de inflexão de renegociação de chaves (\textit{rekeying}).}
\label{fig:ponto_otimo_rekeying}
\end{figure}

\begin{figure}[htbp]
\centering
\includegraphics[width=0.82\textwidth]{dissertacao_03_autonomia_bateria.png}
\caption{Projeção de autonomia operacional de nós sensores alimentados por baterias CR2032, 2x AA e Li-Ion 18650.}
\label{fig:autonomia_bateria}
\end{figure}

\subsection{Discussão da Amortização e Zona Ótima de Rekeying}
\label{subsec:discussao_rekeying}

A inspeção detalhada das Figuras~\ref{fig:protocolo_trend}, \ref{fig:energia_amortizada} e \ref{fig:ponto_otimo_rekeying} revela conclusões de alta relevância para a engenharia de sistemas embarcados:
\begin{enumerate}
    \item \textbf{Redução Drástica do Custo Computacional e Energético}: no regime monomensagem ($N=1$), o custo por pacote é integralmente dominado pelo handshake ($49.681{,}9~\mu\text{s}$ e $24.841{,}0~\mu\text{J}$). Ao estender a sessão para apenas $N=10$ mensagens, o custo unitário cai para $4.972{,}6~\mu\text{s}$ (uma economia de $\mathbf{90\%}$). Para $N=100$ mensagens, o custo despenca para $504{,}4~\mu\text{s}$ ($252{,}2~\mu\text{J}$), atingindo uma \textbf{economia superior a $98{,}9\%$};
    \item \textbf{Determinação do Ponto de Rekeying Ótimo}: a Figura~\ref{fig:ponto_otimo_rekeying} evidencia que a curva de custo atinge seu joelho de saturação assintótica entre $\mathbf{N=50}$ e $\mathbf{N=100}$ pacotes. A partir de $N=100$, os ganhos marginais por mensagem adicional são inferiores a $0{,}1\%$, enquanto o risco criptográfico de exposição de chave acumula-se. Logo, o limiar de \textit{rekeying} ótimo situa-se em $N \in [50, 100]$ mensagens, garantindo \textit{Perfect Forward Secrecy} sem penalizar a vida útil do dispositivo;
    \item \textbf{Impacto na Vida Útil de Baterias Industriais}: a Figura~\ref{fig:autonomia_bateria} comprova que, com sessões amortizadas ($N \ge 100$) em regime de amostragem de 1 mensagem por minuto, uma bateria comum tipo moeda CR2032 (2376 J / 0.66 Wh) oferece autonomia de processamento criptográfico superior a \textbf{12 anos}, demonstrando a plena aplicabilidade do protocolo em nós remotos.
\end{enumerate}

% =============================================================================
% SEÇÃO 4.6: VALIDAÇÃO EXPERIMENTAL DA DEFESA
% =============================================================================
\section{Validação Experimental da Defesa e Inspeção Matemática}
\label{sec:validacao_defesa}

Para consolidar as evidências técnicas perante a banca examinadora, foi desenvolvido e executado o demonstrador interativo da defesa (\texttt{demonstrador\_defesa}), que permitiu simular ataques adversariais controlados e realizar a inspeção matemática detalhada de cada bloco e quadro físico em trânsito.

\subsection{Simulação e Mitigação de Ataque Man-in-the-Middle (MITM)}
\label{subsec:mitm_attack}

O ensaio adversarial consistiu em interceptar o pacote criptográfico em trânsito no meio físico e alterar um único bit ($1\text{ bit flip}$) no texto cifrado ou no \textit{nonce} antes de entregá-lo ao receptor legítimo.

A Figura~\ref{fig:tikz_mitm} ilustra o diagrama de blocos do processo de interceptação e descarte pelo Poly1305.

\begin{figure}[htbp]
\centering
\begin{tikzpicture}[
    node distance=1.5cm and 1.8cm,
    box/.style={draw=black, rounded corners, thick, align=center, fill=blue!8, minimum width=2.6cm, minimum height=1.1cm},
    mitm/.style={draw=red!80!black, fill=red!10, rounded corners, thick, align=center, minimum width=3.2cm, minimum height=1.1cm},
    decision/.style={draw=black, diamond, aspect=2, thick, align=center, fill=yellow!15, minimum width=2.8cm},
    arrow/.style={-Latex, thick}
]
    % Nós
    \node[box] (alice) {\textbf{Emissor (Alice)}\\Gera $C$ e Tag $T_E$};
    \node[mitm, right=of alice] (attacker) {\textbf{Canal / Atacante MITM}\\Adulteração de 1 bit\\$C \to C'$};
    \node[box, right=of attacker] (bob) {\textbf{Receptor (Bob)}\\Recalcula Tag $T_R(C')$};
    \node[decision, below=1.2cm of bob] (check) {$T_E \stackrel{?}{=} T_R$};
    \node[box, fill=green!15, left=2.2cm of check] (accept) {\textbf{Payload Decifrado}\\Mensagem Aceita};
    \node[box, fill=red!20, below=1.0cm of check] (reject) {\textbf{Descarte Imediato}\\Payload Não Revelado\\Alerta de Ataque};

    % Conexões
    \draw[arrow] (alice) -- node[above, font=\footnotesize] {Pacote $(N, C, T_E)$} (attacker);
    \draw[arrow] (attacker) -- node[above, font=\footnotesize] {Pacote Alterado $(N, C', T_E)$} (bob);
    \draw[arrow] (bob) -- (check);
    \draw[arrow] (check) -- node[above, font=\footnotesize] {Sim (Válido)} (accept);
    \draw[arrow] (check) -- node[right, font=\footnotesize] {Não (Falha)} (reject);
\end{tikzpicture}
\caption{Diagrama de fluxo da mitigação e descarte determinístico de ataque MITM pelo Poly1305.}
\label{fig:tikz_mitm}
\end{figure}

Em 100\% dos ensaios realizados (1.000 mensagens com injeção de falhas bit a bit), a tag recalculada divergiu da tag transmitida ($T_E \ne T_R$), resultando na rejeição imediata pelo receptor sem realizar a liberação do payload em memória. A verificação do Poly1305 é realizada em tempo constante, prevenindo vazamentos por oráculos de preenchimento ou temporização.

\subsection{Estrutura do Pacote no Meio Físico (Antena / Cabo)}
\label{subsec:quadro_meio_fisico}

A inspeção matemática permitiu revelar explicitamente a estrutura do pacote transmitido no canal de rádio e os passos de computação simétrica.

A Figura~\ref{fig:tikz_quadro_fisico} detalha a arquitetura do quadro em nível de enlace.

\begin{figure}[htbp]
\centering
\begin{tikzpicture}[
    field/.style={draw=black, thick, align=center, minimum height=1.4cm, fill=cyan!10},
    arrow/.style={-Latex, thick}
]
    % Campos do Pacote
    \node[field, fill=orange!20, minimum width=3.2cm] (nonce) {\textbf{Nonce de Enlace}\\12 Bytes (96 bits)\\Público / Anti-Reúso};
    \node[field, fill=blue!15, minimum width=6.2cm, right=0cm of nonce] (ct) {\textbf{Texto Cifrado (Ciphertext)}\\$L$ Bytes (ex.: 64~B ou 513~B)\\Confidencialidade ($P \oplus K$)};
    \node[field, fill=green!20, minimum width=3.8cm, right=0cm of ct] (tag) {\textbf{Tag Poly1305}\\16 Bytes (128 bits)\\Integridade e Autenticidade};

    % Chaves e Dimensões
    \draw[thick, decoration={brace, mirror, raise=6pt}, decorate]
        (nonce.south west) -- node[below=10pt, font=\small] {\textbf{Overhead Fixo AEAD = 28 Bytes} (12~B Nonce + 16~B Tag)} (tag.south east);

    \draw[thick, decoration={brace, raise=6pt}, decorate]
        (ct.north west) -- node[above=10pt, font=\small] {\textbf{Payload Útil = $L$ Bytes}} (ct.north east);
\end{tikzpicture}
\caption{Formato estrutural do quadro criptográfico AEAD transmitido no meio físico.}
\label{fig:tikz_quadro_fisico}
\end{figure}

\subsection{Inspeção Matemática da Cifragem e Evolução do Nonce}
\label{subsec:inspecao_matematica}

O demonstrador interativo permitiu verificar a conformidade com as regras estritas da RFC 8439:
\begin{enumerate}
    \item \textbf{Evolução do Nonce Anti-Reúso}: o vetor base da sessão $N_0$ de 12 bytes sofre incremento estrito pelo contador de pacote $i$:
    \begin{equation}
        Nonce_i = N_0 + i \quad (i \in [0, 2^{32}-1])
    \end{equation}
    O valor é tratado em ordenação \textit{Little-Endian} nas palavras de 32 bits da matriz de estado do ChaCha20, impedindo colisões de chaves e garantindo que dois pacotes jamais compartilhem o mesmo fluxo de keystream;
    \item \textbf{Cifragem por Operação XOR Bloco a Bloco}: para cada bloco de 64 bytes gerado após as 20 rodadas do ChaCha20, o keystream $K_i$ é combinado diretamente ao payload $P_i$:
    \begin{equation}
        C_i = P_i \oplus K_i
    \end{equation}
    Para o último bloco fracionário (como evidenciado no teste com 513 bytes, onde o bloco 9 possui apenas 1 byte), exatamente 1 byte de keystream (o primeiro byte da primeira palavra em Little-Endian) é consumido, garantindo preservação de banda sem expansão de tamanho de mensagem;
    \item \textbf{Garantia de Fidedignidade na Decifragem}: o receptor verifica a tag Poly1305 sobre o texto cifrado e os dados associados autenticados (AAD). Confirmada a autenticidade, o texto cifrado sofre a mesma operação simétrica de decifragem ($P = C \oplus K$), recuperando fielmente a mensagem original transmitida.
\end{enumerate}

\subsection{Comparativo Estrutural com Padrões Clássicos da Literatura}
\label{subsec:comparativo_literatura}

Para situar a contribuição no estado da arte, a Tabela~\ref{tab:comparativo_estrutural_literatura} sintetiza um comparativo estrutural entre o protocolo proposto e os dois padrões dominantes na literatura e na indústria: RSA-2048 com AES-128 e a suíte clássica NIST P-256 com AES-GCM.

\begin{table}[htbp]
\centering
\caption{Comparação estrutural entre a suíte proposta e os padrões clássicos baseados nas referências da dissertação.}
\label{tab:comparativo_estrutural_literatura}
\begin{tabular}{lccc}
\hline
\textbf{Parâmetro de Avaliação} & \textbf{RSA-2048 + AES-128} & \textbf{NIST P-256 + AES-GCM} & \textbf{Proposta: Ed25519 + X25519 + CC20} \\
\hline
Nível de Segurança & 112 bits & 128 bits & \textbf{128 bits} \\
Tamanho da Chave Pública & 256 bytes & 64 bytes & \textbf{32 bytes ($8\times$ menor que RSA)} \\
Tamanho da Assinatura & 256 bytes & 64 bytes & \textbf{64 bytes ($4\times$ menor que RSA)} \\
Overhead de Handshake em Rede & $\sim 512$ bytes & $\sim 160$ bytes & \textbf{128 bytes (Menor tráfego)} \\
Aceleração de HW no Cortex-A9 & Não & Não & \textbf{Instruções ARX nativas (1 ciclo)} \\
Resistência a Cache-Timing & Vulnerável em software & Vulnerável a S-Box & \textbf{Imune (Sem tabelas de consulta)} \\
Consumo de Energia no Handshake & $\sim 56{,}8\text{ mWh}$ & $\sim 15{,}4\text{ mWh}$ & \textbf{$\sim 12{,}5\text{ mJ}$ ($< 0{,}004\text{ mWh}$)} \\
\hline
\multicolumn{4}{l}{\footnotesize Dados comparativos extraídos das referências da dissertação (Suarez-Albela 2018, Tanksale 2024).} \\
\end{tabular}
\end{table}

A análise da Tabela~\ref{tab:comparativo_estrutural_literatura} evidencia a superioridade da proposta: enquanto o RSA-2048 impõe sobrecargas proibitivas de rede e energia para canais de baixa taxa e o NIST P-256 apresenta riscos de vazamento por canais laterais de \textit{cache-timing} e inversões no corpo $\mathbb{F}_p$, o protocolo híbrido baseado em curvas de Edwards/Montgomery e cifra ChaCha20 atinge segurança de 128 bits com consumo de energia por mensagem na ordem de dezenas de microjoules e chaves compactas de apenas 32 bytes.

% =============================================================================
% SEÇÃO 4.7: CONCLUSÃO DO CAPÍTULO E TRANSIÇÃO
% =============================================================================
\section{Conclusões do Capítulo}
\label{sec:conclusoes_capitulo}

Os resultados experimentais consolidados ao longo de 50 macro-rodadas e 50.000 iterações em ambiente emulado de sistema completo comprovaram com elevado rigor estatístico e determinismo temporal que a combinação de primitivas criptográficas ultraleves de 128 bits (Ed25519, X25519, HKDF-SHA512 e ChaCha20-Poly1305) é plenamente viável para dispositivos embarcados restritos com processadores ARM Cortex-A9.

A caracterização dos quatro eixos operacionais demonstrou que:
\begin{enumerate}
    \item A assinatura e troca de chaves assimétricas possuem comportamento temporal altamente determinístico ($CV < 1{,}8\%$), com a Escada de Montgomery demonstrando escalabilidade assintótica linear e imunidade a canais laterais sob a norma ISO/IEC 17825;
    \item A cifra de fluxo simétrica com autenticação integrada garante vazão superior a 31~MB/s sem vulnerabilidades de \textit{cache-timing}, demandando apenas $4{,}2~\mu\text{s}$ e $2{,}1~\mu\text{J}$ para pacotes típicos de telemetria;
    \item A sobrecarga de enlace fixa de 28 bytes preserva a compatibilidade com os padrões LoRaWAN, Zigbee e BLE sem exigir fragmentação de nível de rede para cargas úteis até 64 bytes;
    \item A amortização por sessão viabiliza uma redução superior a $98{,}9\%$ no consumo de energia e tempo de CPU por mensagem, assegurando vida útil de bateria de mais de uma década em nós sensores.
\end{enumerate}

Com a sustentação empírica e a comprovação matemática de robustez estabelecidas neste capítulo, o próximo capítulo sintetiza as conclusões finais desta dissertação, apresentando as contribuições científicas alcançadas e as perspectivas de trabalhos futuros para implementação em silício e integração com arquiteturas de computação de borda.
